% Propuesta de redacción del capítulo 5 (resultados). Generada por experiments/chapter5.py a partir de
% results/evaluation.json, results/final/specialists.json y results/final/graph_stats.json.
% NO reemplaza secciones/capitulo5.tex: es material para que Ana lo integre.
% Tablas: \begin{table}[H]
\centering
\footnotesize
\renewcommand{\arraystretch}{1.1}
\resizebox{\textwidth}{!}{%
\begin{tabular}{llccccc}
\toprule
\textbf{ID} & \textbf{Sistema} & \textbf{PR-AUC} & \textbf{Macro-F1} & \textbf{Brier} & \textbf{ECE} & \textbf{AURC} \\
\midrule
B0 & Reglas explícitas & 0.555 $\pm$ 0.000 & 0.571 $\pm$ 0.000 & 0.393 $\pm$ 0.000 & 0.014 $\pm$ 0.000 & 0.217 $\pm$ 0.000 \\
 & & {\scriptsize 0.555 [0.549, 0.565]} & {\scriptsize 0.571 [0.571, 0.571]} & {\scriptsize 0.393 [0.389, 0.396]} & {\scriptsize 0.014 [0.013, 0.020]} & {\scriptsize 0.217 [0.201, 0.216]} \\
B1 & Solo texto & 0.735 $\pm$ 0.002 & 0.513 $\pm$ 0.006 & 0.470 $\pm$ 0.000 & 0.002 $\pm$ 0.001 & 0.230 $\pm$ 0.002 \\
 & & {\scriptsize 0.736 [0.721, 0.747]} & {\scriptsize 0.507 [0.488, 0.527]} & {\scriptsize 0.470 [0.463, 0.478]} & {\scriptsize 0.005 [0.000, 0.011]} & {\scriptsize 0.226 [0.216, 0.238]} \\
B2 & Solo daño visual & 0.552 $\pm$ 0.004 & 0.467 $\pm$ 0.015 & 0.502 $\pm$ 0.003 & 0.019 $\pm$ 0.009 & 0.289 $\pm$ 0.007 \\
 & & {\scriptsize 0.553 [0.540, 0.571]} & {\scriptsize 0.476 [0.453, 0.499]} & {\scriptsize 0.500 [0.494, 0.507]} & {\scriptsize 0.019 [0.010, 0.031]} & {\scriptsize 0.285 [0.268, 0.302]} \\
B3 & Solo forense & 0.758 $\pm$ 0.001 & 0.549 $\pm$ 0.005 & 0.472 $\pm$ 0.000 & 0.011 $\pm$ 0.003 & 0.215 $\pm$ 0.003 \\
 & & {\scriptsize 0.758 [0.742, 0.773]} & {\scriptsize 0.554 [0.536, 0.571]} & {\scriptsize 0.472 [0.463, 0.481]} & {\scriptsize 0.011 [0.009, 0.034]} & {\scriptsize 0.214 [0.202, 0.228]} \\
\midrule
B4 & Concatenación + MLP & 0.867 $\pm$ 0.001 & 0.693 $\pm$ 0.010 & 0.264 $\pm$ 0.001 & 0.011 $\pm$ 0.003 & 0.070 $\pm$ 0.001 \\
 & & {\scriptsize 0.868 [0.856, 0.880]} & {\scriptsize 0.689 [0.668, 0.710]} & {\scriptsize 0.264 [0.254, 0.274]} & {\scriptsize 0.010 [0.007, 0.024]} & {\scriptsize 0.070 [0.063, 0.077]} \\
B5 & Fusión tardía calibrada & 0.878 $\pm$ 0.001 & 0.669 $\pm$ 0.013 & 0.264 $\pm$ 0.002 & 0.013 $\pm$ 0.005 & 0.071 $\pm$ 0.002 \\
 & & {\scriptsize 0.878 [0.868, 0.888]} & {\scriptsize 0.664 [0.642, 0.687]} & {\scriptsize 0.263 [0.254, 0.272]} & {\scriptsize 0.011 [0.006, 0.030]} & {\scriptsize 0.071 [0.064, 0.078]} \\
B7 & DBF & 0.850 $\pm$ 0.031 & 0.662 $\pm$ 0.052 & 0.443 $\pm$ 0.122 & 0.142 $\pm$ 0.105 & 0.151 $\pm$ 0.016 \\
 & & {\scriptsize 0.864 [0.852, 0.875]} & {\scriptsize 0.688 [0.658, 0.715]} & {\scriptsize 0.385 [0.378, 0.392]} & {\scriptsize 0.099 [0.083, 0.117]} & {\scriptsize 0.145 [0.133, 0.157]} \\
B8 & MLLM congelado & 0.438 $\pm$ 0.000 & 0.398 $\pm$ 0.000 & 0.600 $\pm$ 0.000 & 0.072 $\pm$ 0.000 & 0.366 $\pm$ 0.000 \\
 & & {\scriptsize 0.438 [0.431, 0.450]} & {\scriptsize 0.398 [0.370, 0.425]} & {\scriptsize 0.600 [0.591, 0.609]} & {\scriptsize 0.072 [0.057, 0.090]} & {\scriptsize 0.366 [0.352, 0.382]} \\
\midrule
\rowcolor{black!6} B6 & DeepSets & 0.869 $\pm$ 0.001 & 0.705 $\pm$ 0.007 & 0.264 $\pm$ 0.001 & 0.010 $\pm$ 0.004 & 0.069 $\pm$ 0.001 \\
\rowcolor{black!6}  & & {\scriptsize 0.870 [0.859, 0.881]} & {\scriptsize 0.711 [0.691, 0.732]} & {\scriptsize 0.263 [0.255, 0.272]} & {\scriptsize 0.012 [0.005, 0.025]} & {\scriptsize 0.068 [0.062, 0.074]} \\
\rowcolor{black!6} G0 & Agregador, $g=1$ & 0.871 $\pm$ 0.002 & 0.698 $\pm$ 0.003 & 0.266 $\pm$ 0.003 & 0.021 $\pm$ 0.004 & 0.071 $\pm$ 0.002 \\
\rowcolor{black!6}  & & {\scriptsize 0.871 [0.861, 0.882]} & {\scriptsize 0.699 [0.675, 0.721]} & {\scriptsize 0.265 [0.256, 0.275]} & {\scriptsize 0.016 [0.012, 0.033]} & {\scriptsize 0.070 [0.064, 0.077]} \\
\rowcolor{black!6} G1 & \modelo{} (agregador con $g$) & 0.851 $\pm$ 0.000 & 0.671 $\pm$ 0.000 & 0.276 $\pm$ 0.000 & 0.015 $\pm$ 0.000 & 0.080 $\pm$ 0.000 \\
\rowcolor{black!6}  & & {\scriptsize 0.851 [0.838, 0.865]} & {\scriptsize 0.671 [0.647, 0.695]} & {\scriptsize 0.276 [0.267, 0.285]} & {\scriptsize 0.015 [0.007, 0.027]} & {\scriptsize 0.080 [0.072, 0.087]} \\
\bottomrule
\end{tabular}}
\caption{Resultados principales en test (evaluado una sola vez). Primera línea: media $\pm$ desviación entre 5 semillas (cada una calibrada con su $T^\ast$); segunda línea: ensamble calibrado e IC 95\% del bootstrap agrupado por \texttt{source\_image\_id} (2\,000 réplicas) sobre el ensamble calibrado. Sin GATv2: G0/G1 son el agregador interpretable sobre el grafo.}
\label{tab:resultados_main}
\end{table}
, tab_hipotesis, tab_ablaciones, tab_oraculo,
% tab_sensibilidad, tab_delta_m, tabla_especialistas, tabla_externa, tabla_grafo.

\section{Cambio de diseño: agregador sin GATv2 y verificación humana}

El sistema evaluado no usa la capa GATv2. La propagación de mensajes aprendida se reemplazó por un
agregador interpretable sobre el mismo grafo de evidencia: para cada afirmación $i$ se calculan
$S_{\text{apoyo}}=\sum_j g_j\,p_{\text{support}}(i,j)$, $S_{\text{contra}}=\sum_j g_j\,p_{\text{contr}}(i,j)$,
$S_{\text{desc}}=\sum_j g_j\,p_{\text{unk}}(i,j)$, la similitud ontológica ponderada y el término forense
$\max_j g_f\,p_{\text{edit}}\,\mathrm{IoU}(\ell_j,\ell_f)$; se agregan por reclamo (máximo y media), junto con conteos,
$\bar c$, la calidad media y la polaridad de las afirmaciones, y un clasificador liviano elegido en validación produce
$p(y\mid x)$. G0 es el mismo agregador con $g=1$ y reemplaza a P1; G1 usa la compuerta y reemplaza a P2. En validación,
el clasificador seleccionado fue logreg para G1 y gboost para G0. Cada decisión se acompaña de un
reporte para revisión humana (foto con máscaras y mapa forense, afirmaciones, grafo con $g$, decisión calibrada,
$U$, $\bar c$, motivo de abstención y las tres piezas de evidencia con mayor peso).

\section{Especialistas}

El especialista narrativo acierta el tipo de daño exacto en 0.500 de los
1129 reclamos de validación con afirmación, 0.689 a nivel de
familia de la ontología HITL, y la polaridad en 0.965; por eso la confiabilidad narrativa $r$ se mide a nivel
de familia. YOLOv8n-seg, ajustado en Kaggle (GPU T4) con el split congelado, obtiene un mAP50 de máscara de
0.118 en validación y 0.096 en test (mAP50-95: 0.055 y
0.041). En CarDD (validación externa, 374 imágenes, solo cajas) el mAP50 de caja es
0.054 con clases mapeadas y 0.077 sin distinguir clase: el detector apenas
transfiere fuera de HITL. TruFor separa las imágenes editadas de las originales en validación con AUROC
0.889 para C3 (copy-move), 0.544 para C4 (inpainting) y 0.774
para C3--C5 en conjunto (C5 usa las mismas imágenes que C3). Ver Tablas~\ref{tab:especialistas} y~\ref{tab:externa}.

\section{Grafos de evidencia}

Los grafos tienen en promedio 5.58 nodos y 13.26
aristas por reclamo en test (6.50 y 16.90 en
entrenamiento). En C0 el 80.0\,\% de las aristas afirmación--daño son de apoyo y el
0.0\,\% de contradicción; en C1 y C5 la contradicción sube a
70.4\,\%. En C2 (sustitución de imagen) la contradicción es
0.1\,\%: el grafo no ve diferencia con C0. La compuerta media de los nodos visuales
es 0.030 en C0, porque la confianza y la calidad regional de las detecciones son bajas; en C9
(desenfoque) el detector no produce nodos visuales y la $g$ de la fuente cae de 0.806 (C0) a
0.007. La $g$ forense media es 0.190 en C0 y 0.364
en C3 (Tabla~\ref{tab:grafo}, Figura de distribución de $g$).

\section{Resultados principales}

La Tabla~\ref{tab:resultados_main} resume test (evaluado una sola vez). G1 obtiene PR-AUC de SUSPICIOUS
0.851 [0.838, 0.865], Macro-F1 0.671, Brier 0.276,
ECE 0.015 y AURC 0.080. G0 obtiene 0.871 [0.861, 0.882],
0.699, 0.265, 0.016 y 0.070; B6 (DeepSets)
0.870, 0.711, 0.263, 0.012 y 0.068.
Ordenados por PR-AUC: B5 (0.878), G0 (0.871), B6 (0.870), B4 (0.868), B7 (0.864), G1 (0.851), B3 (0.758), B1 (0.736), B0 (0.555), B2 (0.553), B8 (0.438). B8 (Qwen2-VL-2B congelado, sin ajuste) queda en 0.438, por debajo
incluso de las reglas B0 (0.555); Macro-F1 0.398 frente a 0.571.
El resultado es negativo para la propuesta: con la evidencia visual disponible, ni el grafo ni la compuerta superan a una
fusión simple.

\section{Contraste de hipótesis}

Con bootstrap pareado agrupado por \texttt{source\_image\_id} (129 imágenes fuente de test, 2\,000 réplicas)
y Holm ($m=3$, $\alpha=0{,}05$) (Tabla~\ref{tab:hipotesis}):

\begin{itemize}
\item \textbf{H1} (G1 frente al mejor sistema sin grafo, elegido en validación: B5; PR-AUC de SUSPICIOUS en C3--C5 frente a
C8--C9, 817 reclamos): diferencia G1$-$B5 = -0.0388 (IC 95\,\% [-0.0530, -0.0246]; $p$ $<$ 0.0005 (0 de 2000 réplicas), $p_{\mathrm{Holm}}$ $<$ 0.0005 (0 de 2000 réplicas)). G1 obtiene 0.854 y B5
0.893. Se rechaza la igualdad, pero a favor de B5: \textbf{H1 no se sostiene}. McNemar sobre
la exactitud global: $b=153$, $c=150$, $p=0.909$.
\item \textbf{H2} (G1 frente a G0 en Brier): Brier(G0)$-$Brier(G1) = -0.0104 (IC 95\,\% [-0.0192, -0.0018]; $p$ 0.0160, $p_{\mathrm{Holm}}$ 0.0160); G1 0.2758, G0
0.2654. La compuerta empeora la calibración: \textbf{H2 no se sostiene} (efecto en contra). McNemar:
$b=87$, $c=128$, $p=0.0062$.
\item \textbf{H3} (G1 frente a A2, sin forense): diferencia de PR-AUC = +0.0618 (IC 95\,\% [+0.0505, +0.0734]; $p$ $<$ 0.0005 (0 de 2000 réplicas), $p_{\mathrm{Holm}}$ $<$ 0.0005 (0 de 2000 réplicas)); G1 0.851, A2
0.789. \textbf{H3 se acepta}: el especialista forense aporta. McNemar: $b=185$, $c=58$.
\end{itemize}

Los valores $p$ reportados como ``$<$ 0.0005'' corresponden a 0 de 2\,000 réplicas con el signo contrario; no son exactamente cero.

\section{Robustez por condición, calibración y abstención}

La exactitud de la decisión final (con abstención contada como INSUFFICIENT) de G1 es 0.674 en C0,
0.845 en C1, 0.000 en C2, 0.651 en C3, 0.008 en C4,
0.953 en C5 y 0.708 en C8 (Tabla~\ref{tab:delta_m}). C2 (sustitución) y C4 (inpainting) son
las condiciones más difíciles para todos los sistemas: $\Delta M_{C2}$ = +0.674 y $\Delta M_{C4}$ = +0.667
para G1; +0.209 y +0.116 para B5. En C8 (transformaciones benignas) G1 acierta
0.708 frente a 0.579 de B5 (fracción decidida SUSPICIOUS en C8: G1
0.000, B5
0.129).
Todos los sistemas calibrados tienen ECE baja (G1 0.015). Los umbrales fijados en validación para una cobertura
de 0.80 dejaron $\tau_c=$ 0.000 (la cobertura de evidencia no discriminó en validación) y $\tau_U=$
0.679; G1 se abstiene en 0.184 de los reclamos etiquetados de test, y solo
0.149 de sus abstenciones caen en C6, C7, C9 o C10: las abstenciones no se concentran
donde deberían, porque esas condiciones ya se resuelven con confianza. Para C10 (OOD), el AUROC de $U$ para separarla del resto
es 0.646 para G1, 0.597 para G0, 0.619 para B6 y 0.571 para B5.

\section{Ablaciones, sensibilidad y oráculo}

Sin forense (A2) la PR-AUC cae a 0.789. Quitar los atributos continuos de las aristas (A4) deja
0.851 (G1: 0.851). Entrenar sin C8 (A5) mantiene la PR-AUC
(0.850) pero hunde el Macro-F1 a 0.460 y la calibración (ECE 0.135).
Sin los 552 reclamos AMBIGUOS (sensibilidad C2--C5), G1 obtiene PR-AUC 0.869
y G0 0.888: la conclusión no cambia (Tabla~\ref{tab:sensibilidad}).
Con percepción oráculo la PR-AUC de G1 sube a 0.866 (narrativa),
0.943 (visual+forense) y 0.965 (ambas); G0 con ambas
llega a 0.970. El cuello de botella es la percepción visual y forense, no el agregador.

\section{Análisis de errores y verificación humana}

Se seleccionaron 20 errores de G1 (sin abstención), repartidos por condición; sus reportes están en
\texttt{review/errores/}. Condiciones de los 20 errores: C0: 1, C1: 4, C2: 4, C3: 4, C4: 4, C5: 3. Tipo de la evidencia con más peso en cada error:
afirmación apoya daño visual (13), forense global (7); pares etiqueta$\to$decisión: SUSPICIOUS$\to$SUPPORTED (19), SUPPORTED$\to$SUSPICIOUS (1). La probabilidad máxima media de esos errores es 0.612. El paquete de
revisión de 50 casos (\texttt{review/}) está listo para que la asesora y la autora lo llenen; no contiene respuestas de
revisión.

\section{Amenazas a la validez}

Los reclamos son escenarios controlados sobre imágenes reales, no expedientes de una aseguradora. El especialista visual es
débil y no transfiere a CarDD. El test no es completamente ciego (se miró al diseñar la QA). C3 y C5 comparten imágenes. La
varianza entre semillas solo cubre al clasificador final: G1 eligió una regresión logística convexa y sus 5 semillas coinciden
(desviación 0). B8 se usó sin ajuste y con un solo prompt.
